\chapter{2018年天津卷（文）}

\section{单选题}

\begin{problem}
设集合 \(A=\{1,2,3,4\}\)，\(B=\{-1,0,2,3\}\)，\(C=\{x\in\R\mid -1\le x<2\}\)，则 \((A\cup B)\cap C=\)（\quad）

\choices
    {\(\{-1,1\}\)}
    {\(\{0,1\}\)}
    {\(\{-1,0,1\}\)}
    {\(\{2,3,4\}\)}

\end{problem}

\begin{answer}
C
\end{answer}

\begin{solution}
由 \(A\cup B=\{-1,0,1,2,3,4\}\)，又 \(C=[-1,2)\)，所以 \((A\cup B)\cap C=\{-1,0,1\}\). 故选 C.
\end{solution}

\begin{problem}
设变量 \(x\)，\(y\) 满足约束条件 \(x+y\le5\)，\(2x-y\le4\)，\(-x+y\le1\)，\(y\ge0\)，则目标函数 \(z=3x+5y\) 的最大值为（\quad）

\choices
    {\(6\)}
    {\(19\)}
    {\(21\)}
    {\(45\)}

\end{problem}

\begin{answer}
C
\end{answer}

\begin{solution}
作出不等式组表示的可行域，其顶点包含 \((-1,0)\)，\((2,0)\)，\((3,2)\)，\((2,3)\). 分别代入 \(z=3x+5y\)，得 \(-3\)，\(6\)，\(19\)，\(21\)，所以目标函数的最大值为 \(21\). 故选 C.
\end{solution}

\begin{problem}
设 \(x\in\R\)，则“\(x^3>8\)”是“\(|x|>2\)”的（\quad）

\choices
    {充分而不必要条件}
    {必要而不充分条件}
    {充要条件}
    {既不充分也不必要条件}

\end{problem}

\begin{answer}
A
\end{answer}

\begin{solution}
由 \(x^3>8\) 得 \(x>2\)，从而 \(|x|>2\). 但 \(|x|>2\) 可得 \(x>2\) 或 \(x<-2\)，不能推出 \(x^3>8\). 故“\(x^3>8\)”是“\(|x|>2\)”的充分而不必要条件.
\end{solution}

\begin{problem}
阅读如图所示的程序框图，运行相应的程序，若输入 \(N\) 的值为 \(20\)，则输出 \(T\) 的值为（\quad）

\texfigure[width=6.0cm,max width=6.0cm]{img_repaint/2018/tianjin_liberal/q4_fig1.tex}

\choices
    {\(1\)}
    {\(2\)}
    {\(3\)}
    {\(4\)}

\end{problem}

\begin{answer}
B
\end{answer}

\begin{solution}
由程序框图知，初始 \(i=2\)，\(T=0\). 当 \(i=2\) 时，\(\frac{20}{2}\) 是整数，故 \(T=1\)，再令 \(i=3\). 当 \(i=3\) 时，\(\frac{20}{3}\) 不是整数，\(T\) 不变，再令 \(i=4\). 当 \(i=4\) 时，\(\frac{20}{4}\) 是整数，故 \(T=2\)，再令 \(i=5\). 此时满足 \(i\ge5\)，输出 \(T=2\). 故选 B.
\end{solution}

\begin{problem}
已知 \(a=\log_3\dfrac72\)，\(b=\left(\dfrac14\right)^{\frac13}\)，\(c=\log_{\frac13}\dfrac15\)，则 \(a\)，\(b\)，\(c\) 的大小关系为（\quad）

\choices
    {\(a>b>c\)}
    {\(b>a>c\)}
    {\(c>b>a\)}
    {\(c>a>b\)}

\end{problem}

\begin{answer}
D
\end{answer}

\begin{solution}
因为 \(1<\frac72<9\)，所以 \(1<a<2\). 又 \(0<b=\sqrt[3]{\frac14}<1\). 因为 \(\left(\frac13\right)^2=\frac19<\frac15<\frac13\)，且函数 \(y=\log_{\frac13}x\) 单调递减，所以 \(1<c<2\). 进一步有 \(3^a=\frac72\)，\(3^c=5\)，故 \(c>a\). 因此 \(c>a>b\). 故选 D.
\end{solution}

\begin{problem}
将函数 \(y=\sin\!\left(2x+\dfrac{\pi}{5}\right)\) 的图象向右平移 \(\dfrac{\pi}{10}\) 个单位长度，所得图象对应的函数（\quad）

\choices
    {在区间 \(\left[-\dfrac{\pi}{4},\dfrac{\pi}{4}\right]\) 上单调递增}
    {在区间 \(\left[-\dfrac{\pi}{4},0\right]\) 上单调递减}
    {在区间 \(\left[\dfrac{\pi}{4},\dfrac{\pi}{2}\right]\) 上单调递增}
    {在区间 \(\left[\dfrac{\pi}{2},\pi\right]\) 上单调递减}

\end{problem}

\begin{answer}
A
\end{answer}

\begin{solution}
平移后的函数为 \(y=\sin\!\left(2\left(x-\frac{\pi}{10}\right)+\frac{\pi}{5}\right)=\sin2x\). 函数 \(y=\sin2x\) 在 \(\left[-\frac{\pi}{4},\frac{\pi}{4}\right]\) 上单调递增. 故选 A.
\end{solution}

\begin{problem}
已知双曲线 \(\dfrac{x^2}{a^2}-\dfrac{y^2}{b^2}=1\)（\(a>0\)，\(b>0\)）的离心率为 \(2\)，过右焦点且垂直于 \(x\) 轴的直线与双曲线交于 \(A\)，\(B\) 两点. 设 \(A\)，\(B\) 到双曲线同一条渐近线的距离分别为 \(d_1\) 和 \(d_2\)，且 \(d_1+d_2=6\)，则双曲线的方程为（\quad）

\choices
    {\(\dfrac{x^2}{3}-\dfrac{y^2}{9}=1\)}
    {\(\dfrac{x^2}{9}-\dfrac{y^2}{3}=1\)}
    {\(\dfrac{x^2}{4}-\dfrac{y^2}{12}=1\)}
    {\(\dfrac{x^2}{12}-\dfrac{y^2}{4}=1\)}

\end{problem}

\begin{answer}
A
\end{answer}

\begin{solution}
设右焦点为 \(F(c,0)\)，则 \(e=\frac ca=2\)，故 \(c=2a\)，\(b^2=c^2-a^2=3a^2\). 渐近线为 \(bx\pm ay=0\). 直线 \(x=c\) 与双曲线交于 A，B，两点到同一条渐近线距离之和为 \(\frac{2bc}{\sqrt{a^2+b^2}}=2b\)，故 \(2b=6\)，得 \(b=3\). 于是 \(a^2=3\)，双曲线方程为 \(\frac{x^2}{3}-\frac{y^2}{9}=1\). 故选 A.
\end{solution}

\begin{problem}
在如图的平面图形中，已知 \(OM=1\)，\(ON=2\)，\(\angle MON=120^\circ\)，\(\overrightarrow{BM}=2\overrightarrow{MA}\)，\(\overrightarrow{CN}=2\overrightarrow{NA}\)，则 \(\overrightarrow{BC}\cdot\overrightarrow{OM}\) 的值为（\quad）

\bitmapfigure[max width=6.0cm]{img/2018/tianjin_liberal/q8_fig1.png}

\choices
    {\(-15\)}
    {\(-9\)}
    {\(-6\)}
    {\(0\)}

\end{problem}

\begin{answer}
C
\end{answer}

\begin{solution}
由 \(\overrightarrow{BM}=2\overrightarrow{MA}\)，得 \(M\) 是 \(BA\) 的三等分点，且 \(\overrightarrow{BA}=3\overrightarrow{MA}\). 由 \(\overrightarrow{CN}=2\overrightarrow{NA}\)，得 \(N\) 是 \(CA\) 的三等分点. 可取点 \(O\) 为原点，记 \(\overrightarrow{OM}=\bs m\)，\(\overrightarrow{ON}=\bs n\)，则 \(|\bs m|=1\)，\(|\bs n|=2\)，\(\bs m\cdot\bs n=1\cdot2\cos120^\circ=-1\). 由分点关系可得 \(\overrightarrow{BC}=3(\bs n-\bs m)\)，所以 \(\overrightarrow{BC}\cdot\overrightarrow{OM}=3(\bs n-\bs m)\cdot\bs m=3(-1-1)=-6\). 故选 C.
\end{solution}
\section{填空题}

\begin{problem}
\(\i\) 是虚数单位，复数 \(\dfrac{6+7\i}{1+2\i}=\fillinblank{}\).

\end{problem}

\begin{answer}
\(4-\i\)
\end{answer}

\begin{solution}
\(\frac{6+7\i}{1+2\i}=\frac{(6+7\i)(1-2\i)}{(1+2\i)(1-2\i)}=\frac{20-5\i}{5}=4-\i\).
\end{solution}

\begin{problem}
已知函数 \(f(x)=\e^x\ln x\)，\(f'(x)\) 为 \(f(x)\) 的导函数，则 \(f'(1)\) 的值为\fillinblank{}.

\end{problem}

\begin{answer}
\(\e\)
\end{answer}

\begin{solution}
由 \(f'(x)=\e^x\ln x+\e^x\cdot\frac1x\)，得 \(f'(1)=\e\).
\end{solution}

\begin{problem}
如图，已知正方体 \(ABCD-A_1B_1C_1D_1\) 的棱长为 \(1\)，则四棱锥 \(A_1-BB_1D_1D\) 的体积为\fillinblank{}.

\bitmapfigure[max width=6.0cm]{img/2018/tianjin_liberal/q11_fig1.png}

\end{problem}

\begin{answer}
\(\dfrac13\)
\end{answer}

\begin{solution}
取平面 \(BB_1D_1D\) 为底面. 在正方体中，四边形 \(BB_1D_1D\) 为矩形，\(|BB_1|=1\)，\(|BD|=\sqrt2\)，故底面积为 \(\sqrt2\). 点 \(A_1\) 到平面 \(BB_1D_1D\) 的距离等于点 A 到直线 \(BD\) 的距离，为 \(\frac{\sqrt2}{2}\). 所以四棱锥体积为 \(\frac13\cdot\sqrt2\cdot\frac{\sqrt2}{2}=\frac13\).
\end{solution}

\begin{problem}
在平面直角坐标系中，经过三点 \((0,0)\)，\((1,1)\)，\((2,0)\) 的圆的方程为\fillinblank{}.

\end{problem}

\begin{answer}
\(x^2+y^2-2x=0\)
\end{answer}

\begin{solution}
设圆的方程为 \(x^2+y^2+Dx+Ey+F=0\). 代入 \((0,0)\)，\((1,1)\)，\((2,0)\)，得 \(F=0\)，\(2+D+E=0\)，\(4+2D=0\). 解得 \(D=-2\)，\(E=0\)，故圆的方程为 \(x^2+y^2-2x=0\).
\end{solution}

\begin{problem}
已知 \(a\)，\(b\in\R\)，且 \(a-3b+6=0\)，则 \(2^a+\dfrac1{8^b}\) 的最小值为\fillinblank{}.

\end{problem}

\begin{answer}
\(\dfrac14\)
\end{answer}

\begin{solution}
由 \(a-3b+6=0\) 得 \(a-3b=-6\). 又 \(2^a+\frac1{8^b}=2^a+2^{-3b}\ge2\sqrt{2^a\cdot2^{-3b}}=2\sqrt{2^{-6}}=\frac14\). 当 \(2^a=2^{-3b}\) 时等号成立，故最小值为 \(\frac14\).
\end{solution}

\begin{problem}
已知 \(a\in\R\)，函数 \(f(x)= \begin{cases}
x^2+2x+a-2, & x\le0,\\
-x^2+2x-2a, & x>0.
\end{cases}\)
若对任意 \(x\in[-3,+\infty)\)，\(f(x)\le|x|\) 恒成立，则 \(a\) 的取值范围是\fillinblank{}.

\end{problem}

\begin{answer}
\(\left[\dfrac18,2\right]\)
\end{answer}

\begin{solution}
当 \(-3\le x\le0\) 时，\(|x|=-x\)，由 \(x^2+2x+a-2\le -x\) 得 \(a\le -x^2-3x+2\). 函数 \(-x^2-3x+2\) 在 \([-3,0]\) 上的最小值为 \(2\)，故 \(a\le2\). 当 \(x>0\) 时，\(|x|=x\)，由 \(-x^2+2x-2a\le x\) 得 \(a\ge\frac{-x^2+x}{2}\). 函数 \(\frac{-x^2+x}{2}\) 在 \((0,+\infty)\) 上的最大值为 \(\frac18\)，故 \(a\ge\frac18\). 因此 \(a\in\left[\frac18,2\right]\).
\end{solution}
\section{解答题}

\begin{problem}

已知某校甲，乙，丙三个年级的学生志愿者人数分别为 \(240,160,160\). 现采用分层抽样的方法从中抽取 \(7\) 名同学去某敬老院参加献爱心活动.

\begin{enumerate}
\item 应从甲，乙，丙三个年级的学生志愿者中分别抽取多少人？
\item 设抽出的 \(7\) 名同学分别用 A，B，C，D，E，\(F\)，\(G\) 表示，现从中随机抽取 \(2\) 名同学承担敬老院的卫生工作.
\begin{enumerate}
\item 试用所给字母列举出所有可能的抽取结果；
\item 设 \(M\) 为事件“抽取的 \(2\) 名同学来自同一年级”，求事件 \(M\) 发生的概率.
\end{enumerate}
\end{enumerate}

\end{problem}

\begin{solution}
（1）因为甲，乙，丙三个年级的学生志愿者人数之比为 \(240:160:160=3:2:2\)，所以采用分层抽样从中抽取 \(7\) 名同学，应从甲，乙，丙三个年级的学生志愿者中分别抽取 \(3\) 人，\(2\) 人，\(2\) 人.

（2）（i）从抽出的 \(7\) 名同学中随机抽取 \(2\) 名同学的所有可能结果为 \(\{A,B\}\)，\(\{A,C\}\)，\(\{A,D\}\)，\(\{A,E\}\)，\(\{A,F\}\)，\(\{A,G\}\)，\(\{B,C\}\)，\(\{B,D\}\)，\(\{B,E\}\)，\(\{B,F\}\)，\(\{B,G\}\)，\(\{C,D\}\)，\(\{C,E\}\)，\(\{C,F\}\)，\(\{C,G\}\)，\(\{D,E\}\)，\(\{D,F\}\)，\(\{D,G\}\)，\(\{E,F\}\)，\(\{E,G\}\)，\(\{F,G\}\)，共 \(21\) 种.

（2）不妨设来自甲年级的是 A，B，C，来自乙年级的是 D，E，来自丙年级的是 \(F\)，\(G\). 则抽取的 \(2\) 名同学来自同一年级的所有可能结果为 \(\{A,B\}\)，\(\{A,C\}\)，\(\{B,C\}\)，\(\{D,E\}\)，\(\{F,G\}\)，共 \(5\) 种. 所以 \(P(M)=\frac5{21}\).
\end{solution}

\begin{problem}

在 \(\triangle ABC\) 中，内角 \(A\)，\(B\)，\(C\) 所对的边分别为 \(a\)，\(b\)，\(c\). 已知 \(b\sin A=a\cos\!\left(B-\dfrac{\pi}{6}\right)\).

\begin{enumerate}
\item 求角 \(B\) 的大小；
\item 设 \(a=2\)，\(c=3\)，求 \(b\) 和 \(\sin(2A-B)\) 的值.
\end{enumerate}

\end{problem}

\begin{solution}
（1）由正弦定理得 \(b\sin A=a\sin B\). 因为 \(a>0\)，所以 \(\sin B=\cos\!\left(B-\frac{\pi}{6}\right)=\sin\!\left(\frac{2\pi}{3}-B\right)\). 又 \(0<B<\pi\)，得 \(B=\frac{\pi}{3}\).

（2）由余弦定理得 \(b^2=a^2+c^2-2ac\cos B=4+9-2\cdot2\cdot3\cdot\frac12=7\)，故 \(b=\sqrt7\). 由 \(b\sin A=a\cos\!\left(B-\frac{\pi}{6}\right)\)，得 \(\sin A=\frac2{\sqrt7}\cdot\frac{\sqrt3}{2}=\sqrt{\frac37}\). 因为 \(a<c\)，所以 \(A\) 为锐角，故 \(\cos A=\frac2{\sqrt7}\). 于是 \(\sin2A=\frac{4\sqrt3}{7}\)，\(\cos2A=\frac17\)，从而 \(\sin(2A-B)=\sin2A\cos B-\cos2A\sin B=\frac{3\sqrt3}{14}\).
\end{solution}

\begin{problem}

如图，在四面体 \(ABCD\) 中，\(\triangle ABC\) 是等边三角形，平面 \(ABC\perp\) 平面 \(ABD\)，点 \(M\) 为棱 \(AB\) 的中点，\(AB=2\)，\(AD=2\sqrt3\)，\(\angle BAD=90^\circ\).

\begin{enumerate}
\item 求证：\(AD\perp BC\)；
\item 求异面直线 \(BC\) 与 \(MD\) 所成角的余弦值；
\item 求直线 \(CD\) 与平面 \(ABD\) 所成角的正弦值.
\end{enumerate}

\bitmapfigure[max width=6.0cm]{img/2018/tianjin_liberal/q17_fig1.png}

\end{problem}

\begin{solution}
（1）因为平面 \(ABC\perp\) 平面 \(ABD\)，平面 \(ABC\cap\) 平面 \(ABD=AB\)，且 \(AD\perp AB\)，\(AD\subset\) 平面 \(ABD\)，所以 \(AD\perp\) 平面 \(ABC\). 又 \(BC\subset\) 平面 \(ABC\)，故 \(AD\perp BC\).

（2）取棱 \(AC\) 的中点 \(N\)，连接 \(MN\)，\(ND\). 因为 \(M\) 为棱 \(AB\) 的中点，所以 \(MN//BC\)，故 \(\angle DMN\) 或其补角为异面直线 \(BC\) 与 \(MD\) 所成的角. 在直角三角形 \(DAM\) 中，\(AM=1\)，\(DM=\sqrt{AD^2+AM^2}=\sqrt{13}\). 又 \(AN=1\)，且 \(AD\perp AC\)，所以 \(DN=\sqrt{AD^2+AN^2}=\sqrt{13}\)，且 \(MN=1\). 在等腰三角形 \(DMN\) 中，\(\cos\angle DMN=\frac{DM^2+MN^2-DN^2}{2DM\cdot MN}=\frac1{2\sqrt{13}}=\frac{\sqrt{13}}{26}\). 所以异面直线 \(BC\) 与 \(MD\) 所成角的余弦值为 \(\frac{\sqrt{13}}{26}\).

（3）连接 \(CM\). 因为 \(\triangle ABC\) 为等边三角形，\(M\) 为 \(AB\) 的中点，所以 \(CM\perp AB\)，\(CM=\sqrt3\). 又平面 \(ABC\perp\) 平面 \(ABD\)，且 \(CM\subset\) 平面 \(ABC\)，所以 \(CM\perp\) 平面 \(ABD\). 因此 \(\angle CDM\) 为直线 \(CD\) 与平面 \(ABD\) 所成的角. 在直角三角形 \(CAD\) 中，\(CD=\sqrt{AC^2+AD^2}=4\). 所以 \(\sin\angle CDM=\frac{CM}{CD}=\frac{\sqrt3}{4}\).
\end{solution}

\begin{problem}

设 \(\{a_n\}\) 是等差数列，其前 \(n\) 项和为 \(S_n\)（\(n\in\N^*\)）；\(\{b_n\}\) 是等比数列，公比大于 \(0\)，其前 \(n\) 项和为 \(T_n\)（\(n\in\N^*\)）. 已知 \(b_1=1\)，\(b_3=b_2+2\)，\(b_4=a_3+a_5\)，\(b_5=a_4+2a_6\).

\begin{enumerate}
\item 求 \(S_n\) 和 \(T_n\)；
\item 若 \(S_n+(T_1+T_2+\cdots+T_n)=a_n+4b_n\)，求正整数 \(n\) 的值.
\end{enumerate}

\end{problem}

\begin{solution}
（1）设等比数列 \(\{b_n\}\) 的公比为 \(q\). 由 \(b_1=1\)，\(b_3=b_2+2\)，得 \(q^2=q+2\). 因为 \(q>0\)，所以 \(q=2\)，故 \(b_n=2^{n-1}\)，\(T_n=2^n-1\). 设等差数列 \(\{a_n\}\) 的首项为 \(a_1\)，公差为 \(d\). 由 \(b_4=a_3+a_5\)，得 \(a_1+3d=4\). 由 \(b_5=a_4+2a_6\)，得 \(3a_1+13d=16\). 解得 \(a_1=1\)，\(d=1\)，故 \(a_n=n\)，\(S_n=\frac{n(n+1)}2\).

（2）由（1）知 \(T_1+T_2+\cdots+T_n=(2^1+2^2+\cdots+2^n)-n=2^{n+1}-n-2\). 代入 \(S_n+(T_1+T_2+\cdots+T_n)=a_n+4b_n\)，得 \(\frac{n(n+1)}2+2^{n+1}-n-2=n+2^{n+1}\). 整理得 \(n^2-3n-4=0\)，解得 \(n=4\) 或 \(n=-1\). 因为 \(n\in\N^*\)，所以 \(n=4\).
\end{solution}

\begin{problem}

设椭圆 \(\dfrac{x^2}{a^2}+\dfrac{y^2}{b^2}=1\)（\(a>b>0\)）的右顶点为 \(A\)，上顶点为 \(B\). 已知椭圆的离心率为 \(\dfrac{\sqrt5}{3}\)，\(|AB|=\sqrt{13}\).

\begin{enumerate}
\item 求椭圆的方程；
\item 设直线 \(l:y=kx\)（\(k<0\)）与椭圆交于 \(P\)，\(Q\) 两点，\(l\) 与直线 \(AB\) 交于点 \(M\)，且点 \(P\)，\(M\) 均在第四象限. 若 \(\triangle BPM\) 的面积是 \(\triangle BPQ\) 面积的 \(2\) 倍，求 \(k\) 的值.
\end{enumerate}

\end{problem}

\begin{solution}
（1）设椭圆的焦距为 \(2c\)，由 \(e=\frac ca=\frac{\sqrt5}{3}\)，得 \(\frac{c^2}{a^2}=\frac59\). 又 \(a^2=b^2+c^2\)，所以 \(2a=3b\). 因为 \(|AB|=\sqrt{a^2+b^2}=\sqrt{13}\)，所以 \(a=3\)，\(b=2\). 故椭圆方程为 \(\frac{x^2}{9}+\frac{y^2}{4}=1\).

（2）设点 \(P(x_1,y_1)\)，点 \(M(x_2,y_2)\)，则点 \(Q(-x_1,-y_1)\). 由面积关系可得 \(|PM|=2|PQ|\)，从而 \(x_2-x_1=2(x_1-(-x_1))\)，即 \(x_2=5x_1\). 直线 \(AB\) 的方程为 \(2x+3y=6\). 由 \(y=kx\) 与 \(2x+3y=6\)，得 \(x_2=\frac6{3k+2}\). 由 \(y=kx\) 与椭圆方程，得 \(x_1=\frac6{\sqrt{9k^2+4}}\). 由 \(x_2=5x_1\)，得 \(\sqrt{9k^2+4}=5(3k+2)\). 两边平方整理得 \(18k^2+25k+8=0\)，解得 \(k=-\frac89\) 或 \(k=-\frac12\). 当 \(k=-\frac89\) 时 \(x_2<0\)，不合题意；当 \(k=-\frac12\) 时符合题意. 故 \(k=-\frac12\).
\end{solution}

\begin{problem}

设函数 \(f(x)=(x-t_1)(x-t_2)(x-t_3)\)，其中 \(t_1\)，\(t_2\)，\(t_3\in\R\)，且 \(t_1\)，\(t_2\)，\(t_3\) 是公差为 \(d\) 的等差数列.

\begin{enumerate}
\item 若 \(t_2=0\)，\(d=1\)，求曲线 \(y=f(x)\) 在点 \((0,f(0))\) 处的切线方程；
\item 若 \(d=3\)，求 \(f(x)\) 的极值；
\item 若曲线 \(y=f(x)\) 与直线 \(y=-(x-t_2)-6\sqrt3\) 有三个互异的公共点，求 \(d\) 的取值范围.
\end{enumerate}

\end{problem}

\begin{solution}
因为 \(t_1\)，\(t_2\)，\(t_3\) 是公差为 \(d\) 的等差数列，所以 \(t_1=t_2-d\)，\(t_3=t_2+d\)，从而 \(f(x)=(x-t_2+d)(x-t_2)(x-t_2-d)\).

（1）当 \(t_2=0\)，\(d=1\) 时，\(f(x)=x(x-1)(x+1)=x^3-x\)，故 \(f(0)=0\)，\(f'(0)=-1\). 所求切线方程为 \(y=-x\)，即 \(x+y=0\).

（2）当 \(d=3\) 时，\(f(x)=(x-t_2)^3-9(x-t_2)\). 所以 \(f'(x)=3(x-t_2)^2-9\). 令 \(f'(x)=0\)，得 \(x=t_2-\sqrt3\) 或 \(x=t_2+\sqrt3\). 由导数符号变化知，\(f(x)\) 在 \(x=t_2-\sqrt3\) 处取得极大值 \(6\sqrt3\)，在 \(x=t_2+\sqrt3\) 处取得极小值 \(-6\sqrt3\).

（3）令 \(u=x-t_2\). 曲线 \(y=f(x)\) 与直线 \(y=-(x-t_2)-6\sqrt3\) 有三个互异公共点，等价于方程 \(u^3+(1-d^2)u+6\sqrt3=0\) 有三个互异实根. 设 \(g(u)=u^3+(1-d^2)u+6\sqrt3\)，则 \(g'(u)=3u^2+1-d^2\). 若 \(d^2\le1\)，则 \(g'(u)\ge0\)，不可能有三个互异实根. 若 \(d^2>1\)，令 \(u_1=-\sqrt{\frac{d^2-1}{3}}\)，\(u_2=\sqrt{\frac{d^2-1}{3}}\). 方程有三个互异实根等价于 \(g(u_1)>0\) 且 \(g(u_2)<0\). 其中 \(g(u_1)>0\) 恒成立，而 \(g(u_2)<0\) 等价于 \(\frac{2\sqrt3}{9}(d^2-1)^{\frac32}>6\sqrt3\)，即 \((d^2-1)^{\frac32}>27\)，从而 \(d^2>10\). 因此 \(d<-\sqrt{10}\) 或 \(d>\sqrt{10}\).
\end{solution}
